% ===================================================================
%  TABELL Nr 2 — Människokroppen. — Rasten. — Lekarna.
%  Traduction 2026 d'apres texto/io, controlee sur texto/fr et
%  texto/es. Consigne : voir texto/sv/10-tabell-01.tex.
% ===================================================================

%%K t02-apar-1 apar
\VUsaut{4.0mm}
\VUtitre{15.0pt}{60}{TABELL Nr 2}
\VUsaut{2.0mm}
\VUtitre{11.4pt}{0}{Människokroppen. --- Rasten.}
\VUtitre{11.4pt}{0}{Lekarna.}

%%K t02-tit-0 sub
\VUtitre{13.2pt}{0}{Människokroppen.}
\VUsaut{2.4mm}
\VUcentre{10.2pt}{0}{\textit{(Enkel lektion i naturvetenskap.)}}
\VUsaut{3.0mm}

%%K t02-tit-1 sub pos=t02-01-1
\VUpk{10.2pt}{40}{Huvudet.}
\VUsaut{3.0mm}

%%K t02-01-1 p
1. --- Människokroppens förnämsta delar är : \VUgras{huvudet}, \VUgras{bålen} och
\VUgras{lemmarna}.

%%K t02-02-1 p
2. --- Huvudet omfattar två delar : \VUgras{skallen}, som innehåller
\VUgras{hjärnan} och som \VUgras{håret} täcker, och \VUgras{ansiktet}.

%%K t02-03-1 p
3. --- På vår teckning ser vi varken skallen eller hjärnan; men vi ser mycket tydligt
ansiktets viktiga delar.

%%K t02-04-1 p
4. --- Här är först \VUgras{pannan}, som tyder på ett kraftigt förstånd, ett ständigt verksamt
tänkande. \VUgras{Tinningarna} är mycket fria. \VUgras{Ögat} är livligt och vidöppet. Ett tätt
\VUgras{ögonbryn} och långa
\VUgras{ögonfransar} skyddar det.

%%K t02-05-1 p
5. --- Här är dessutom \VUgras{näsan} och \VUgras{näsborrarna}. Näsan tjänar oss för att lukta
och för att andas. På var sida om näsan är \VUgras{kinderna}, längre bort \VUgras{öronen}.
Ansiktets nedre del består av \VUgras{munnen} och
\VUgras{hakan}.

%%K t02-06-1 p
6. --- Vi ser dem inte särskilt väl, för \VUgras{mustaschen} och
\VUgras{polisongerna} döljer dem för oss. Men vi vet, att munnen har de två
\VUgras{läpparna}, \VUgras{överkäken} och \VUgras{underkäken} med deras
\VUgras{tänder}, och \VUgras{tungan}. Med munnen talar och äter vi. Med tungan
smakar vi födan.

%%K t02-07-1 p
7. --- Ögat, örat, näsan och munnen är, liksom fingrarna, de fem sinnenas organ :
\VUgras{synen}, \VUgras{hörseln}, \VUgras{luktsinnet}, \VUgras{smaken},
\VUgras{känseln}.

%%K t02-08-1 p
8. --- Huvudet är alltså en av människokroppens mest intressanta delar.

%%K t02-tit-2 sub
\VUsaut{4.4mm}
\VUpk{10.2pt}{40}{Bålen.}
\VUsaut{3.0mm}

%%K t02-09-1 p
9. --- \VUgras{Halsen} förenar huvudet med bålen. Den omfattar \VUgras{strupen} och
\VUgras{nacken}.

%%K t02-10-1 p
10. --- Bålen omfattar också många väsentliga organ. I \VUgras{bröstet} finns
\VUgras{lungorna}, med vilka vi andas, och \VUgras{hjärtat} som är livets
medelpunkt. När hjärtat upphör att slå, dör den levande varelsen.

%%K t02-11-1 p
11. --- \VUgras{Revbenen} håller upp bröstet.

%%K t02-12-1 p
12. --- Bålens övriga delar är \VUgras{sidan}, \VUgras{ryggen},
\VUgras{skuldrorna} och \VUgras{buken}. Vi lägger oss på sidan eller på ryggen
för att sova, vi bär bördorna på vår skuldra.

%%K t02-13-1 p
13. --- I buken finns \VUgras{magen} och \VUgras{tarmarna}. De tjänar oss för att smälta
födan.

%%K t02-tit-3 sub
\VUsaut{4.4mm}
\VUpk{10.2pt}{40}{Lemmarna.}
\VUsaut{3.0mm}

%%K t02-14-1 p
14. --- Vi har fyra \VUgras{lemmar} : två \VUgras{armar} och två \VUgras{ben}. Med armarna
tar, håller, lyfter och samlar vi föremålen; med benen står, går och springer vi.

%%K t02-15-1 p
15. --- \VUgras{Skuldran} förenar armen med kroppen. En mycket kraftig
\VUgras{muskel}, bicepsen, rör armen som \VUgras{armbågen} förenar med
\VUgras{underarmen}. \VUgras{Handloven} utgör \VUgras{handens} led, som slutar
med \VUgras{fingrarna} och \VUgras{naglarna}. På handen urskiljer vi en
\VUgras{ven} (och en artär) i vilken \VUgras{blodet} flyter.

%%K t02-16-1 p
16. --- Benets delar är : \VUgras{låret}, \VUgras{knäet} och \VUgras{knävecket},
\VUgras{vaden}, vars \VUgras{kött} täcks av \VUgras{huden}, \VUgras{fotknölen}
och \VUgras{foten}. Benets \VUgras{ben} är mycket grova och mycket starka. Vid fotens ände är
\VUgras{tårna}. De övriga delarna är \VUgras{fotsulan} och
\VUgras{hälen}.

%%K t02-17-1 p
17. --- Sådan är den nästan fullständiga beskrivningen av människokroppen.

%%K t02-17-2 p
\textit{Anm.} --- \textit{Med hjälp av uttrycket : « jag har ont i... (huvudet o. s. v.) » kan
läraren använda hela denna ordskatt på nytt ur synpunkten av det goda eller dåliga}
\VUgras{kroppstillståndet.}

%%K t02-tit-4 sub
\VUsaut{5.0mm}
\VUpk{10.2pt}{40}{Gymnastik.}
\VUsaut{2.4mm}
\VUcentre{10.2pt}{0}{\textit{(Berättelse av en av eleverna.)}}
\VUsaut{3.0mm}

%%K t02-18-1 p
18. --- För att vara smidiga, viga och friska måste vi öva våra ben och våra armar. Därför
leker vi och utövar \VUgras{gymnastik}.

%%K t02-19-1 p
19. --- Under veckan äger dessa två övningar rum på läroverkets gård (se tabell nr 1). Men på
torsdagar och söndagar går vi till en stor \VUgras{gräsplan}, där vi har utrymme att springa
och roa oss.

%%K t02-20-1 p
20. --- Våra lärare och våra föräldrar följer oss ibland dit; men en
\VUgras{gymnastiklärare} \textsuperscript{(12)} leder ju övningarna.

%%K t02-21-1 p
21. --- På gräsplanen, till vänster om ingången, står
\VUgras{gymnastikredskapen} \textsuperscript{(1)}, vi klättrar uppför
\VUgras{stegen} \textsuperscript{(2)}, vi hänger upp och ner i \VUgras{ringarna}
\textsuperscript{(3)} och i \VUgras{trapetsen} \textsuperscript{(4)}, vi hissar oss med
händerna upp till toppen av \VUgras{repstegen} \textsuperscript{(5)} eller av
\VUgras{knutrepet} \textsuperscript{(6)}.

%%K t02-22-1 p
22. --- Under tiden hoppar våra kamrater över \VUgras{trähästen} \textsuperscript{(7)} eller
\VUgras{barren} \textsuperscript{(11)}. Andra
\VUgras{boxas} \textsuperscript{(8)} eller \VUgras{fäktar} \textsuperscript{(9)}
med mask och med \VUgras{florett}. Slutligen \VUgras{lyfter} andra
\VUgras{skivstänger} \textsuperscript{(10)}.

%%K t02-tit-5 sub
\VUsaut{4.6mm}
\VUpk{10.2pt}{40}{Lekarna.}
\VUsaut{3.0mm}

%%K t02-23-1 p
23. --- Runt gymnasterna leker pojkar och flickor olika \VUgras{lekar}. Flickorna roar sig med
sitt \VUgras{tunnband} \textsuperscript{(14)}; de hoppar
\VUgras{rep} \textsuperscript{(13)} eller bjuder sina bröder och kusiner till ett
parti \VUgras{krocket} \textsuperscript{(15)}. De har nöje av att stöta bort motståndarens
\VUgras{klot} med sin \VUgras{träklubba}, just i det ögonblick då han tror sig lätt gå igenom
bågen!

%%K t02-24-1 p
24. --- Pojkarna föredrar de lekar som kräver mer muskler. Gärna leker de
\VUgras{kägelspel} \textsuperscript{(16)} där de skickligt slår ner käglorna med
ett \VUgras{klot} \textsuperscript{(17)}. De försmår inte heller att leka med
\VUgras{snurra} \textsuperscript{(18)} och \VUgras{bilboquet}
\textsuperscript{(19)} eller till och med \VUgras{grodspel} \textsuperscript{(20)}. Men deras
favoritlekar är ju \VUgras{kulorna} \textsuperscript{(21)} och \VUgras{bollen mot muren}
\textsuperscript{(22)}, för \VUgras{växthusets} \textsuperscript{(26)} mur är mycket lämplig
för att låta den lätta bollen studsa tillbaka.

%%K t02-25-1 p
25. --- Lille Johan, som går på sina \VUgras{styltor} \textsuperscript{(24)}, är mycket
skicklig och kan mycket väl hålla balansen. Bredvid honom leker några kamrater
\VUgras{kurragömma} \textsuperscript{(25)} i det växthuset och bakom
\VUgras{häcken}. Andra föredrar \VUgras{slaggissning} \textsuperscript{(28)}
eller \VUgras{fjäderbollen} \textsuperscript{(29)} som flyger från det ena
\VUgras{racketet} till det andra.

%%K t02-noto-1 noto
\VUnotes{143.2mm}{%
(*) Barn- eller pojklekarna har ofta rent nationella namn. Vi har handlat för att benämna dem
på Ido så gott som möjligt, antingen genom att låta oss inspireras av själva handlingen,
varigenom de kännetecknas, eller genom att välja det nationella namnet i det ena eller andra
landet, när det inte är alltför oriktigt. T. ex. : \VUgras{tunnleken} (på tyska och franska)
är
\VUgras{grodleken} \textsuperscript{(20)} (på engelska) i vilken deltagarna
försöker kasta sina runda brickor genom munnen på den metallgroda som med gapande mun står på
den genomborrade lådan (tunnan).
\VUgras{Skulderhoppet} \textsuperscript{(30)} skiljer sig från
\VUgras{ryggahoppet} bara i detta : för att hoppa över huvudet stöder
deltagarna händerna på kamratens skuldror, medan de vid « \VUgras{ryggahoppet} » stöder
händerna bara på hans rygg. --- Eller också är några av deltagarna böjda medan de andra ska
hoppa över deras rygg med händerna stödda på skuldran; de första måste tåla stöten utan att
böja sig ihop, de andra måste hoppa utan att förlora balansen (Se själva tabellen).%
}

%%K t02-26-1 p
26. --- Mitt på gräsplanen står två träd. Vid foten av ett av de två har sju eller åtta pojkar
ordnat ett parti \VUgras{skulderhopp} \textsuperscript{(30)} (*). Inte i alla länder känner
man den leken; men alla vet vad
\VUgras{ryggahoppet} är, som pojkarna inte leker denna gång.

%%K t02-tit-6 sub
\VUsaut{5.0mm}
\VUpk{10.2pt}{40}{Lekarna i fria luften.}
\VUsaut{3.4mm}

%%K t02-27-1 p
27. \VUgras{Blindbocken} \textsuperscript{(31)} är också mycket rolig; flickor och pojkar
sätter i tur och ordning \VUgras{ögonbindeln} för ögonen och griper de klumpiga som låter sig
fångas.

%%K t02-28-1 p
28. --- Mitt på gräsplanen, här är en mycket fransk men något våldsam lek : det är
\VUgras{björnhonan} \textsuperscript{(32)} som är mycket mer bullersam än den så lugna leken
med \VUgras{drake} \textsuperscript{(33)} eller till och med än den engelska leken
\VUgras{tennis} \textsuperscript{(34)}.

%%K t02-29-1 p
29. --- Den sistnämnda sporten är de stora elevernas glädje. Våra lärare själva utövar den
gärna. Den kräver stor skicklighet med vighet och den behagar mig mycket. Jag överlåter åt
flickorna lekarna med \VUgras{gungan} \textsuperscript{(35)}.

%%K t02-30-1 p
30. --- Jag tycker inte om en annan engelsk lek som kanske skulle bli farlig.

%%K t02-30-1b p
Jag menar \VUgras{fotbollen} \textsuperscript{(36)} som gläder min äldre bror. Jag föredrar
vår gamla \VUgras{bomlek} \textsuperscript{(38)} som är lika hälsosam och verkligen mindre rå.

%%K t02-tit-7 sub
\VUsaut{4.6mm}
\VUpk{10.2pt}{40}{Cykling och bollspel.}
\VUsaut{3.0mm}

%%K t02-31-1 p
31. --- Också gärna åker jag runt gräsplanen på \VUgras{cykel} \textsuperscript{(37)}.

%%K t02-32-1 p
32. --- Under nästa år ska jag lära mig \VUgras{ridningen} \textsuperscript{(39)}. Så vacker
är en häst som graciöst skrider eller travar, och som i galopp hoppar över hindren!

%%K t02-33-1 p
33. --- På sommaren lär vi oss \VUgras{simkonsten} \textsuperscript{(40)}. Vi dyker med
förtjusning och badar i flodens klara och friska vatten. Ibland släpper vi slutligen upp en
\VUgras{ballong} \textsuperscript{(41)} av papper som majestätiskt stiger upp i luften, så
snart den är fylld med varm luft.

%%K t02-34-1 p
34. --- Det finns också många andra lekar, men jag skulle aldrig sluta att räkna upp dem och
man inser av denna sammanfattning, att man på torsdagarna inte har särskilt tråkigt på vår
vackra gräsplan.

%%K t02-34-2 p
Anm. --- \textit{Läraren kan lätt fullständiga det kapitlet; genom att utförligare beskriva
var och en av de viktigaste lekarna.}
\VUsaut{6.0mm}
\VUfilet{22mm}
